\scanpage{13}\mspage{12}
ici la notion de raffinement \uline{stricte}. Dans \unsure{GP|GR} \uline{V} p. 19,20, il faut supposer, pour $F\overset{s}{\ll}\unsure{G|C}$, dans a) \deleted{que} [en plus de $\forall\,A\in F$, $G^A$ a un plus petit élément] que pour $A\in F$, $B\in G$, $A\cap B$ soit ou bien vide, ou une \deleted{\unsure{face de}} sous-strate de $A$ :
\deleted{En termes de la famille induite par $G$ sur $|F|$ (où on ne retient que les traces sur $|F|$ des $B\in G$ telles que $\mathring{B}\cap|F|\neq\unsure{\phi|\varnothing}$).} \deleted{cela signifie ceci} \deleted{si on veut que} \deleted{Cela signifie \unsure{aussi}}

Il faut faire \unsure{une|un} \unsure{pareille|pareil} \deleted{\unsure{strict}} \unsure{durcissement|changement|opération} pour la notion de compatibilité (la renforçant en ``compatibilité stricte'', ou ``polyédrale'') -- qui signifie que, \unsure{non seulement|bien seulement} $F\cup G$ est une \deleted{\unsure{famille}} \inserted{figure \unsure{ensembliste}}, admettant $F$ et $G$ comme sous-figure, mais que c'est une figure ensembliste \uline{polyédrale}.

En résumé, \deleted{pour} la formulation des axiomes spéciaux, il faut distinguer ceux relatifs aux ``ateliers polyédraux'', et ceux relatifs aux ``ateliers \unsure{spéciaux|spacieux}''.

Quand on en viendra \unsure{aux} plongements des $\mathrm{Fig}(\unsure{L|\mathcal{L}})$, dans le premier cas il faudra poser (pour les ateliers \unsure{polyédraux} ensemblistes, à ne pas confondre avec ateliers ensemblistes polyédraux, qui n'existent \unsure{plus} !)
% slanted
\margin{propriété qui \unsure{remplace|figure} \unsure{Axiome|Axiom} \unsure{des|du} \unsure{Cat PL|Cat P|Ch.} \unsure{VI|II} p.52}
que
\[
\deleted{\unsure{F}}\unsure{\mathrm{mulens}|\mathrm{mulont}|\mathrm{mulons}}(F)
\overset{\mathrm{pol}}{\ll}\mathrm{mulens}\,G\Longrightarrow F\ll G.
\]
